\documentclass[11pt]{article}
% The preamble every project's paper shares: % The preamble every project's paper shares: % The preamble every project's paper shares: \input{../../preamble} from
% projects/<project>/paper/. Boards come from chessboard and xskak, plots from
% pgfplots; every number a paper prints comes from its generated.tex.
\usepackage[LSBC4,T1]{fontenc}
\usepackage[utf8]{inputenc}
\usepackage{lmodern}
\usepackage[margin=1in]{geometry}
\usepackage{microtype}
\usepackage{amsmath}
\usepackage{booktabs}
\usepackage{tabularx}
\usepackage{array}
\usepackage{xcolor}
\usepackage[obeyspaces,spaces]{url}
\usepackage{xskak}
\usepackage{pgfplots}
\pgfplotsset{compat=1.18}
\usepackage{caption}
\usepackage{subcaption}
\usepackage{enumitem}
\usepackage[section]{placeins}
\usepackage[hidelinks]{hyperref}

\setchessboard{boardfontencoding=LSBC4, boardfontfamily=skaknew, blackfieldcolor=black!22,
  whitefieldcolor=white, setfontcolors, showmover=false, label=true, labelfontsize=6pt,
  boardfontsize=14pt}

% A move or line in UCI, and a position in FEN, both breakable at a slash.
\newcommand{\uci}[1]{\texttt{#1}}
\newcommand{\fen}[1]{{\footnotesize\path{#1}}}
\newcolumntype{L}{>{\raggedright\arraybackslash}X}

\setlength{\parskip}{0.4em}
\setlength{\parindent}{0pt}
 from
% projects/<project>/paper/. Boards come from chessboard and xskak, plots from
% pgfplots; every number a paper prints comes from its generated.tex.
\usepackage[LSBC4,T1]{fontenc}
\usepackage[utf8]{inputenc}
\usepackage{lmodern}
\usepackage[margin=1in]{geometry}
\usepackage{microtype}
\usepackage{amsmath}
\usepackage{booktabs}
\usepackage{tabularx}
\usepackage{array}
\usepackage{xcolor}
\usepackage[obeyspaces,spaces]{url}
\usepackage{xskak}
\usepackage{pgfplots}
\pgfplotsset{compat=1.18}
\usepackage{caption}
\usepackage{subcaption}
\usepackage{enumitem}
\usepackage[section]{placeins}
\usepackage[hidelinks]{hyperref}

\setchessboard{boardfontencoding=LSBC4, boardfontfamily=skaknew, blackfieldcolor=black!22,
  whitefieldcolor=white, setfontcolors, showmover=false, label=true, labelfontsize=6pt,
  boardfontsize=14pt}

% A move or line in UCI, and a position in FEN, both breakable at a slash.
\newcommand{\uci}[1]{\texttt{#1}}
\newcommand{\fen}[1]{{\footnotesize\path{#1}}}
\newcolumntype{L}{>{\raggedright\arraybackslash}X}

\setlength{\parskip}{0.4em}
\setlength{\parindent}{0pt}
 from
% projects/<project>/paper/. Boards come from chessboard and xskak, plots from
% pgfplots; every number a paper prints comes from its generated.tex.
\usepackage[LSBC4,T1]{fontenc}
\usepackage[utf8]{inputenc}
\usepackage{lmodern}
\usepackage[margin=1in]{geometry}
\usepackage{microtype}
\usepackage{amsmath}
\usepackage{booktabs}
\usepackage{tabularx}
\usepackage{array}
\usepackage{xcolor}
\usepackage[obeyspaces,spaces]{url}
\usepackage{xskak}
\usepackage{pgfplots}
\pgfplotsset{compat=1.18}
\usepackage{caption}
\usepackage{subcaption}
\usepackage{enumitem}
\usepackage[section]{placeins}
\usepackage[hidelinks]{hyperref}

\setchessboard{boardfontencoding=LSBC4, boardfontfamily=skaknew, blackfieldcolor=black!22,
  whitefieldcolor=white, setfontcolors, showmover=false, label=true, labelfontsize=6pt,
  boardfontsize=14pt}

% A move or line in UCI, and a position in FEN, both breakable at a slash.
\newcommand{\uci}[1]{\texttt{#1}}
\newcommand{\fen}[1]{{\footnotesize\path{#1}}}
\newcolumntype{L}{>{\raggedright\arraybackslash}X}

\setlength{\parskip}{0.4em}
\setlength{\parindent}{0pt}

% Written by projects/lesson-videos/paper/generate.mjs; do not edit by hand.

\newcommand{\BaseCommit}{c98822d}
\newcommand{\CandCommit}{4d0e9d3}
\newcommand{\ChangeCommit}{f7a8b12}
\newcommand{\Days}{twelve}
\newcommand{\DaysCap}{Twelve}
\newcommand{\WriterRuns}{24}
\newcommand{\FirstDay}{October 13}
\newcommand{\LastDay}{October 28}
\newcommand{\TrainN}{six}
\newcommand{\HeldN}{six}
\newcommand{\TrainDays}{Oct 13, Oct 15, Oct 18, Oct 20, Oct 24 and Oct 27}
\newcommand{\HeldDays}{Oct 14, Oct 17, Oct 19, Oct 22, Oct 25 and Oct 28}
\newcommand{\GamesMin}{one}
\newcommand{\GamesMax}{14}
\newcommand{\GamesAll}{62}
\newcommand{\WordsPerSecond}{2.6}
\newcommand{\MinDecoyWords}{40}
\newcommand{\Spent}{123}
\newcommand{\Cap}{216}
\newcommand{\Budget}{180}
\newcommand{\WriterSpend}{87}
\newcommand{\GraderSpend}{36}
\newcommand{\SevenDayFirst}{13}
\newcommand{\SevenDayLast}{19}
\newcommand{\BasePassed}{twelve}
\newcommand{\BaseRounds}{2.3}
\newcommand{\BaseHeard}{8.2}
\newcommand{\BaseWords}{1084}
\newcommand{\BaseScenes}{16.3}
\newcommand{\BaseBoardScenes}{11.8}
\newcommand{\BaseQuestions}{9.6}
\newcommand{\BaseTeach}{1.7}
\newcommand{\BaseCost}{3.32}
\newcommand{\BaseMinutes}{14}
\newcommand{\BaseVerdictLines}{4.5}
\newcommand{\BaseVerdictMin}{three}
\newcommand{\BaseVerdictMax}{six}
\newcommand{\BaseLinked}{13}
\newcommand{\BaseLinkedN}{11}
\newcommand{\BaseMomentsN}{88}
\newcommand{\CandPassed}{twelve}
\newcommand{\CandRounds}{2.3}
\newcommand{\CandHeard}{8.6}
\newcommand{\CandWords}{1135}
\newcommand{\CandScenes}{16.3}
\newcommand{\CandBoardScenes}{11.6}
\newcommand{\CandQuestions}{9.9}
\newcommand{\CandTeach}{2.1}
\newcommand{\CandCost}{3.94}
\newcommand{\CandMinutes}{17}
\newcommand{\CandVerdictLines}{0.0}
\newcommand{\CandVerdictMin}{zero}
\newcommand{\CandVerdictMax}{zero}
\newcommand{\CandLinked}{43}
\newcommand{\CandLinkedN}{37}
\newcommand{\CandMomentsN}{87}
\newcommand{\LengthCap}{1.3}
\newcommand{\LengthRatio}{1.07 {\footnotesize (0.97 to 1.17)}}
\newcommand{\TrainTruth}{+0.83 {\footnotesize (+0.50 to +1.00)}}
\newcommand{\TrainCoaching}{+1.00 {\footnotesize (+1.00 to +1.00)}}
\newcommand{\TrainListening}{+1.00 {\footnotesize (+1.00 to +1.00)}}
\newcommand{\TrainOverall}{+1.00 {\footnotesize (+1.00 to +1.00)}}
\newcommand{\HeldTruth}{+0.67 {\footnotesize (+0.33 to +1.00)}}
\newcommand{\HeldCoaching}{+1.00 {\footnotesize (+1.00 to +1.00)}}
\newcommand{\HeldListening}{+0.67 {\footnotesize (+0.33 to +1.00)}}
\newcommand{\HeldOverall}{+1.00 {\footnotesize (+1.00 to +1.00)}}
\newcommand{\AllTruth}{+0.75 {\footnotesize (+0.50 to +1.00)}}
\newcommand{\AllCoaching}{+1.00 {\footnotesize (+1.00 to +1.00)}}
\newcommand{\AllListening}{+0.83 {\footnotesize (+0.58 to +1.00)}}
\newcommand{\AllOverall}{+1.00 {\footnotesize (+1.00 to +1.00)}}
\newcommand{\TruthWins}{nine}
\newcommand{\TruthLosses}{zero}
\newcommand{\TruthTies}{three}
\newcommand{\CoachingWins}{twelve}
\newcommand{\CoachingLosses}{zero}
\newcommand{\CoachingTies}{zero}
\newcommand{\ListeningWins}{ten}
\newcommand{\ListeningLosses}{zero}
\newcommand{\ListeningTies}{two}
\newcommand{\OverallWins}{twelve}
\newcommand{\OverallLosses}{zero}
\newcommand{\OverallTies}{zero}
\newcommand{\Decided}{92}
\newcommand{\SecondShare}{53}
\newcommand{\SameBroken}{zero}
\newcommand{\DecoyCoaching}{zero}
\newcommand{\DecoyLostListening}{twelve}
\newcommand{\DecoyLostOverall}{twelve}
\newcommand{\DecoyAdded}{115}
\newcommand{\BaseAboutFirst}{53}
\newcommand{\BasePurposeNamed}{56}
\newcommand{\BaseTakeaway}{68}
\newcommand{\BaseLessonNamed}{14}
\newcommand{\BaseLessonFits}{75}
\newcommand{\BaseBetterSound}{100}
\newcommand{\BaseEngineRecital}{7}
\newcommand{\BaseThenWhat}{68}
\newcommand{\BaseHabitNamed}{49}
\newcommand{\BaseReplacement}{58}
\newcommand{\BaseOpenQuestionAll}{47}
\newcommand{\BasePraiseNamed}{87}
\newcommand{\BaseMoments}{4.8}
\newcommand{\BaseVerdictScenes}{4.9}
\newcommand{\BaseFalse}{1.17}
\newcommand{\BaseUnsupported}{1.92}
\newcommand{\BaseMisleading}{2.25}
\newcommand{\BaseBoardFalse}{1.67}
\newcommand{\BaseAsked}{62}
\newcommand{\CandAboutFirst}{70}
\newcommand{\CandPurposeNamed}{55}
\newcommand{\CandTakeaway}{77}
\newcommand{\CandLessonNamed}{43}
\newcommand{\CandLessonFits}{74}
\newcommand{\CandBetterSound}{100}
\newcommand{\CandEngineRecital}{9}
\newcommand{\CandThenWhat}{66}
\newcommand{\CandHabitNamed}{51}
\newcommand{\CandReplacement}{60}
\newcommand{\CandOpenQuestionAll}{54}
\newcommand{\CandPraiseNamed}{92}
\newcommand{\CandMoments}{4.4}
\newcommand{\CandVerdictScenes}{0.3}
\newcommand{\CandFalse}{0.42}
\newcommand{\CandUnsupported}{2.00}
\newcommand{\CandMisleading}{3.17}
\newcommand{\CandBoardFalse}{1.08}
\newcommand{\CandAsked}{69}
\newcommand{\ItemsTable}{%
\begin{tabular}{lrr}
\toprule
In the moment scenes & Today's code & The change \\
\midrule
Says what the position is about before the move & 53\% {\footnotesize (57)} & 70\% {\footnotesize (53)} \\
Says what the better move does on the board & 56\% {\footnotesize (57)} & 55\% {\footnotesize (53)} \\
Gives an idea to reuse & 68\% {\footnotesize (57)} & 77\% {\footnotesize (53)} \\
\midrule
Names a study lesson & 14\% {\footnotesize (57)} & 43\% {\footnotesize (53)} \\
\quad and the lesson teaches what the scene teaches & 75\% {\footnotesize (8)} & 74\% {\footnotesize (23)} \\
The better move is sound by the evidence & 100\% {\footnotesize (34)} & 100\% {\footnotesize (32)} \\
Recites an engine line without saying what it does & 7\% {\footnotesize (57)} & 9\% {\footnotesize (53)} \\
Follows the better move to their best reply & 68\% {\footnotesize (34)} & 66\% {\footnotesize (32)} \\
Names the habit behind his move & 49\% {\footnotesize (57)} & 51\% {\footnotesize (53)} \\
Ends on the habit that replaces it & 58\% {\footnotesize (57)} & 60\% {\footnotesize (53)} \\
\midrule
Any board scene: asks his move openly first & 47\% {\footnotesize (96)} & 54\% {\footnotesize (97)} \\
Praise names what the good move saw & 87\% {\footnotesize (15)} & 92\% {\footnotesize (12)} \\
\bottomrule
\end{tabular}
}
\newcommand{\FalseDiff}{$-$0.75 {\footnotesize ($-$1.42 to $-$0.08)}}
\newcommand{\BoardDiff}{$-$0.58 {\footnotesize ($-$1.67 to +0.50)}}
\newcommand{\AskedDiff}{+7 {\footnotesize ($-$5 to +17)}}
\newcommand{\UnsoundDiff}{0 {\footnotesize (0 to 0)}}
\newcommand{\VerdictDiff}{$-$4.5 {\footnotesize ($-$5.1 to $-$4.0)}}
\newcommand{\RuleTable}{%
\begin{tabular}{lrl}
\toprule
Condition & Estimate (95\% interval) & Holds \\
\midrule
Coaching on the train days above 0 & +1.00 {\footnotesize (+1.00 to +1.00)} & yes \\
Coaching on the held-out days above 0 & +1.00 {\footnotesize (+1.00 to +1.00)} & yes \\
Overall not below 0 & +1.00 {\footnotesize (+1.00 to +1.00)} & yes \\
False claims per script not up & $-$0.75 {\footnotesize ($-$1.42 to $-$0.08)} & yes \\
False board claims per script not up & $-$0.58 {\footnotesize ($-$1.67 to +0.50)} & yes \\
The check passes as often & 12 of 12 {\footnotesize (today's: 12)} & yes \\
Controls clean & same: 0 broken; decoy: 0 coaching wins & yes \\
Verdict lines per script down & $-$4.5 {\footnotesize ($-$5.1 to $-$4.0)} & yes \\
\quad and the asked moves made more often, points & +7 {\footnotesize ($-$5 to +17)} & \textbf{no} \\
At most 1.3 times as long & 1.07 {\footnotesize (0.97 to 1.17)} & yes \\
Better moves as sound, points & 0 {\footnotesize (0 to 0)} & yes \\
\bottomrule
\end{tabular}
}
\newcommand{\Kept}{no}
\newcommand{\DayTable}{%
\begin{tabular}{llrccccrrrrr}
\toprule
 & & & \multicolumn{4}{c}{Judge, from the change's side} & Length & \multicolumn{2}{c}{Verdict lines} & \multicolumn{2}{c}{Asked moves} \\
\cmidrule(lr){4-7}\cmidrule(lr){9-10}\cmidrule(lr){11-12}
Day & Split & Games & Truth & Coaching & Listening & Overall & ratio & Today & Change & Today & Change \\
\midrule
Oct 13 & train & 14 & + & + & + & + & 0.80 & 5 & 0 & 100 & 61 \\
Oct 14 & held out & 6 & + & + & + & + & 0.88 & 5 & 0 & 54 & 89 \\
Oct 15 & train & 8 & + & + & + & + & 1.23 & 6 & 0 & 50 & 67 \\
Oct 17 & held out & 5 & + & + & + & + & 0.94 & 3 & 0 & 67 & 75 \\
Oct 18 & train & 1 & 0 & + & + & + & 1.03 & 3 & 0 & 67 & 50 \\
Oct 19 & held out & 2 & + & + & + & + & 1.33 & 4 & 0 & 100 & 100 \\
Oct 20 & train & 9 & + & + & + & + & 1.38 & 5 & 0 & 42 & 75 \\
Oct 22 & held out & 9 & 0 & + & + & + & 1.05 & 4 & 0 & 53 & 73 \\
Oct 24 & train & 1 & + & + & + & + & 1.08 & 5 & 0 & 42 & 50 \\
Oct 25 & held out & 1 & 0 & + & 0 & + & 1.08 & 4 & 0 & 58 & 58 \\
Oct 27 & train & 4 & + & + & + & + & 1.09 & 4 & 0 & 61 & 60 \\
Oct 28 & held out & 2 & + & + & 0 & + & 0.92 & 6 & 0 & 50 & 67 \\
\bottomrule
\end{tabular}
}


\title{Does the daily review teach better when it teaches the move a lesson teaches?\\[0.3em]
\large The lesson pick, its prompt and a check on verdict talk, replayed on \Days{} past days}
\author{Chess Coach research lab}
\date{}

\begin{document}
\maketitle

\begin{abstract}
The app's daily routine writes a review of the player's games every day, with a narrated video.
He asked it to stop telling him what the engine thinks of a position and to teach the way people
learn: at each key moment, a sound move a player at his level can find, what it does on the board,
and the idea behind it from the study lessons he follows. One change does all three parts: the
review's \texttt{teach} command picks the move a lesson teaches among the moves the engine says
hold for him, the routine's prompt asks for that move, what it does and the lesson's rule, and the
routine's checker refuses a video line that says who is better. We replayed the whole review on
\Days{} past days, once with today's code and once with the change, with the routine's own model,
tools and checker. Every review passed the check. The change's videos spoke no verdict line on any
day, against \BaseVerdictMin{} to \BaseVerdictMax{} a day today. Judging blind in both orders, a
model preferred the change's coaching and its whole script on all \Days{} days, and both controls
held. An audit against the evidence found fewer false claims in the change's scripts
(\CandFalse{} per script against \BaseFalse), with more misleading ones (\CandMisleading{} against
\BaseMisleading), mostly lesson rules and findability worded loosely. Its moment scenes said what
the position is about before the move more often (\CandAboutFirst\% against \BaseAboutFirst\%), but
what the better move does no more often (\CandPurposeNamed\% against \BasePurposeNamed\%), so the
share of the three moves he asked for rose by \AskedDiff{} points, an interval that does not clear
zero. That is the one condition of the rule fixed before the first run that did not hold, so by its
rule the change is not kept. A second round, designed after reading the first round's scene audits
and run on the same days, added one sentence to the prompt: the scene that shows his move ends on
the better move and what it does. With it every moment scene said what the better move does, the
share of the asked moves rose by \TwoAskedDiff{} points, the judge preferred the change on
\TwoCoachingWins{} of the \Days{} days, and every condition of the rule held, so the second round is
kept. Its scripts carried more false claims than today's (\TwoFalse{} per script against
\BaseFalse), inside the noise. The owner,
watching a day replayed with it, found that it still names a lesson and plays a line without saying
why each move works; the app's next version, which reads every move of both lines on the board, is
not measured here.
\end{abstract}

\section{The question}

Each day the routine has a model review the day's games and write the annotations the app shows
and a narrated video script: scenes with a board from one of his games and beats with a spoken
line, a caption, the moves played on the board and pauses where he is meant to think. Before the
change, the better move it taught at a turning point was the move players at his band most often
choose among those the engine says hold, else the engine's move, and its videos often said who
stood better: ``level'', ``you are winning'', ``the engine agrees''.

The player asked for the review to teach from his study lessons instead: a sound move he can find,
what it does on the board and the idea behind it, from a lesson when one fits, even when it is not
the engine's first choice, as long as the engine says it keeps his position. The lesson-moves
study measured the move such a review would teach on every turning point of his archive: a lesson
teaches it at most turning points, it gives up a fraction of a win\% against the engine's first
line, and players at his band choose it more often than that first line. This study asks whether
the review the routine writes with the change teaches better, without becoming less true or much
longer.

\section{The change}

The change is the app at \texttt{\ChangeCommit}, in three parts that ship together:
\begin{itemize}[nosep]
\item \textbf{The pick.} \texttt{coach review teach} reads what each move that holds for him is for
  (it meets a threat, wins material, trades queens, breaks with a pawn, puts a rook on an open file,
  and so on) and the lesson chapters that teach it, and teaches the move his band plays most often
  among those a lesson teaches for a purpose his move lacked, with the chapters and their rules.
\item \textbf{The prompt.} The daily prompt asks for that move, what it does in teach's words and
  the chapter's rule in the writer's own, a moment that opens with what the position is about, and
  no line that says who is better.
\item \textbf{The check.} \texttt{daily check} refuses a video line that says who is better, by the
  same rule this study counts verdict lines with. It applies from the day after the change ships;
  the replayed days are older, so the change's cases run an eval-only commit
  (\texttt{\CandCommit}) that moves that date before them and changes nothing else.
\end{itemize}
Today's code is the app's main at \texttt{\BaseCommit}.

\section{The cases and the writer}

A case is a past day whose review landed on the app's main with a video script: \Days{} days from
\FirstDay{} to \LastDay, with \GamesMin{} to \GamesMax{} games each (\GamesAll{} in all). Each config
has its own case of each day: a checkout of the app at the config's code with the coach data as
main held it just before that day's review landed and the day's archived engine analysis, with the
day's brief written again by that code's \texttt{coach daily prepare}, offline. The review itself,
its annotations and its video script, is what a writer run writes. The split was fixed before any
run, interleaved by date: train is \TrainDays; held out is \HeldDays.

The writer runs as the routine does: Opus 5.5, headless, with nothing inherited from the session
that launched it, the study lessons' server attached as the live routine has it, and a short
preamble that names the day, says which steps ran and asks for the review, done when the routine's
check passes. Each config wrote every day once (\WriterRuns{} runs in all).

\section{The graders}

A script with no model reads each run: the final check, the check rounds, files changed, commands
it should not have run, how often it ran \texttt{teach}, the video lines that say who is better by
the change's own rule (run on both configs' scripts), the annotations' moments that link a lesson,
the script's shape, its length heard once (spoken words at \WordsPerSecond{} a second plus every
pause), cost and minutes.

Sonnet 5.5 audits each video script alone against the evidence it was written from: the brief
without its bulk, the day's games and turning points, the engine digest, how often players at his
band find the engine's move, and for every board scene the position as the archive, the engine,
Maia-2 and the board code read it, with the engine's best reply after each move the game did not
play; and the study lessons, every lesson's title and what it teaches, with the sections of each
lesson a script names. It lists the claims the evidence contradicts (\emph{false}), does not carry
(\emph{unsupported}) or slants (\emph{misleading}), and for each board scene its kind, its false
board claims, whether it says who is better, and for a moment scene whether it says what the
position is about before the move, what the better move does, an idea to reuse, a lesson and
whether the lesson fits, whether the better move is sound, and the coaching moves of the earlier
video study.

Sonnet 5.5 also judges each day's two scripts blind, in both orders, on truth, coaching, listening
and overall, with the player's ask in its instructions; a side wins a criterion only when both
orders agree. Two controls check the judge on every day: today's script against itself, and against
a decoy that adds as many spoken words as the change added that day (at least \MinDecoyWords), in
beats that restate their scene with nothing new.

\section{The rule}

A day scores $+1$ when the change wins a criterion, $-1$ when it loses and $0$ for a tie; intervals
are 95\% percentile bootstrap intervals over days. The change ships only when every condition of
Table~\ref{tab:rule} holds, fixed before its first run.

\begin{table}[h]
\centering
\small\RuleTable
\caption{The rule, fixed before the change's first run. Coaching and overall are the judge's scores
per day; the claim rows, the asked moves and the unsound better moves are the audit's numbers per
script, the change minus today's code; the verdict lines are the script's own count; the length is
the change's script heard once over today's on the same day. Kept: \Kept.}
\label{tab:rule}
\end{table}

\section{What the judge preferred}

The judge preferred the change's coaching on all \Days{} days, and its whole script overall on all
\Days: \TrainCoaching{} on the train days and \HeldCoaching{} held out. Truth went to the change on
\TruthWins{} days and to today's code on \TruthLosses, with \TruthTies{} ties; listening to the change
on \ListeningWins{} and to today's code on \ListeningLosses. Table~\ref{tab:days} gives each day. Of
the \Decided{} decided verdicts in single orders, \SecondShare\% went to the script shown second.

The controls held. Today's script against itself tied every criterion on every day
(\SameBroken{} broken). The decoy, today's script with about \DecoyAdded{} spoken words of
restatement added, won coaching on \DecoyCoaching{} days and lost listening on
\DecoyLostListening{} and overall on \DecoyLostOverall: the judge does not reward added words.

\section{What the writing did}

\begin{table}[h]
\centering
\small\ItemsTable
\caption{What the audit finds in the moment scenes, pooled over the \Days{} days: the share of the
scenes that owe an item and make it, with the number owing it in brackets.}
\label{tab:items}
\end{table}

\paragraph{No verdict talk.} Today's code spoke \BaseVerdictLines{} lines a day that say who is
better (\BaseVerdictMin{} to \BaseVerdictMax); the change spoke none. The audit, reading every board
scene, found \BaseVerdictScenes{} scenes a day saying who is better in today's scripts and
\CandVerdictScenes{} in the change's.

\paragraph{The asked moves.} The change's moment scenes said what the position is about before
the move in \CandAboutFirst\% of them (today's: \BaseAboutFirst\%) and gave an idea to reuse in
\CandTakeaway\% (\BaseTakeaway\%), but said what the better move does on the board in
\CandPurposeNamed\% (\BasePurposeNamed\%). Per day, the share of the three it made rose by
\AskedDiff{} points (Table~\ref{tab:items}). Every moment scene owes all three, including the scenes
that show what he played before the better move, which is part of why neither config makes them in
every scene.

\paragraph{Lessons.} The change's annotations link a study lesson on \CandLinked\% of their moments
(\CandLinkedN{} of \CandMomentsN), today's on \BaseLinked\% (\BaseLinkedN{} of \BaseMomentsN). The audit
found the change's moment scenes naming a lesson in \CandLessonNamed\% of them, and the lesson
fitting what the scene teaches in \CandLessonFits\% of those.

\paragraph{Truth.} The change's scripts carried \CandFalse{} false claims per script against
\BaseFalse{} today, a difference of \FalseDiff, and \CandBoardFalse{} false board claims against
\BaseBoardFalse{} (\BoardDiff). They carried more misleading claims, \CandMisleading{} per script
against \BaseMisleading, and about as many unsupported ones (\CandUnsupported{} against
\BaseUnsupported). Most of the added misleading claims are of two kinds: a lesson's rule said in
words the lesson does not use, and a line on how often his band finds a move followed by one that
reads as if he had played it. The better moves both
taught were sound by the evidence in every moment scene that owed one.

\paragraph{Length, the check and cost.} Heard once, the change's videos run \CandHeard{} minutes on
average against \BaseHeard{} (a ratio of \LengthRatio{} on the same day). Every run passed the final
check, in \CandRounds{} check rounds on average against \BaseRounds. The change ran \texttt{teach}
\CandTeach{} times a day (today's code: \BaseTeach). A run cost \CandCost{} dollars against \BaseCost{}
and took \CandMinutes{} minutes against \BaseMinutes. The eval spent \Spent{} dollars of a \Budget{}
budget (\WriterSpend{} on writing, \GraderSpend{} on grading), and the subscription's seven-day window
went from \SevenDayFirst\% to \SevenDayLast\%.

\begin{table}[h]
\centering
\scriptsize\setlength{\tabcolsep}{4pt}\DayTable
\caption{Each day: the judge's verdict on each criterion from the change's side ($+$ the change, $-$
today's code, 0 a tie), the change's length over today's, the verdict lines each spoke, and the share
of the asked moves each made.}
\label{tab:days}
\end{table}

\section{Round 2: the scene of his move names the better move}

Round 1's audit marks, for each moment scene, whether it teaches a better move and whether it says
what that move does. Split that way (Table~\ref{tab:scenes}), the gap was structural: in both
configs the scenes that teach a better move nearly always said what it does, and the scenes that
show only his move never did, though the audit holds every moment scene to all three asked moves.
Round 2 changes one thing, a sentence in the routine's prompt: a segment that plays his move in
one scene and the better move in the next says what the better move is for in both, the scene of
his move ending on the better move and what it does, in \texttt{teach}'s words, before the next
scene plays it. The change is the app at \texttt{\TwoChangeCommit}, run as \texttt{\TwoCommit}
with the verdict check moved onto the replayed days as in round 1. Everything else was fixed before
the run and is round 1's: the \Days{} days and their split, the writer, the graders and the rule.
Today's code was not run again: round 2 reads round 1's baseline runs and their audits unchanged.
The decoy was made again from today's script to the length of round 2's.

\begin{table}[h]
\centering
\small\ScenesTable
\caption{Moment scenes that say what the better move does, by whether the audit reads the scene as
teaching a better move or as showing only his move, pooled over the \Days{} days, with the number of
scenes in brackets. In round 2 every scene named a better move, so the audit read none as showing
only his move.}
\label{tab:scenes}
\end{table}

\paragraph{The rule.} Every condition held (Table~\ref{tab:rule2}), so round 2 is kept. The judge
preferred its coaching on \TwoCoachingWins{} days and tied on \TwoCoachingTies, and its whole script
on \TwoOverallWins; truth went to round 2 on \TwoTruthWins{} days, to today's code on
\TwoTruthLosses{} and tied on \TwoTruthTies{} (Table~\ref{tab:days2}). The decoy, today's script with
about \TwoDecoyAdded{} words of restatement added, won no day's coaching and lost listening on all
\Days. Every run passed the final check.

\begin{table}[h]
\centering
\small\TwoRuleTable
\caption{Round 2 against the rule fixed before round 1, over the \Days{} days unless a row says
otherwise. The judge's rows are the mean verdict from round 2's side ($+1$ round 2, $-1$ today's
code), the audit's rows per script, round 2 minus today's code.}
\label{tab:rule2}
\end{table}

\paragraph{The asked moves.} Round 2's moment scenes said what the better move does in
\TwoPurposeNamed\% of them (today's: \BasePurposeNamed\%; round 1: \CandPurposeNamed\%), said what the
position is about first in \TwoAboutFirst\% and gave an idea to reuse in \TwoTakeaway\%
(Table~\ref{tab:items2}). Per day the share of the three rose by \TwoAskedDiff{} points, an
interval clear of zero, where round 1's was \AskedDiff.

\paragraph{What it cost.} Round 2's scripts carried \TwoFalse{} false claims per script against
\BaseFalse{} today (\TwoFalseDiff) and round 1's \CandFalse, \TwoBoardFalse{} false board claims
against \BaseBoardFalse{} (\TwoBoardDiff), and \TwoMisleading{} misleading ones against
\BaseMisleading: more than today on each, inside the noise, where round 1 had fewer false claims
than today. They followed the better move to their best reply in \TwoThenWhat\% of the scenes
that owed it, against \BaseThenWhat\% today, but in about as many scenes (\TwoThenMade{} of
\TwoThenOwed, against \BaseThenMade{} of \BaseThenOwed): the scenes of his move now name a better
move, so the audit holds them to following it too, and they leave that to the next scene. Its videos ran
\TwoHeard{} minutes heard once, \TwoLengthRatio{} times today's on the same day, and spoke no verdict
line. A run cost \TwoCost{} dollars; the round spent \TwoSpent{} dollars and both rounds
\BothSpent{}, under the \TwoCap{} the owner approved for the second round.

\begin{table}[h]
\centering
\small\ItemsTwoTable
\caption{What the audit finds in the moment scenes in each round, pooled over the \Days{} days: the
share of the scenes that owe an item and make it, with the number owing it in brackets.}
\label{tab:items2}
\end{table}

\begin{table}[h]
\centering
\scriptsize\setlength{\tabcolsep}{4pt}\TwoDayTable
\caption{Each day in round 2: the judge's verdict on each criterion from round 2's side ($+$ round 2,
$-$ today's code, 0 a tie), the length over today's, the verdict lines it spoke, and the share of
the asked moves each made.}
\label{tab:days2}
\end{table}

\section{What the app can take}

\begin{itemize}[nosep]
\item By the rule, round 1 is not kept and round 2 is: with the sentence that ends the scene of his
  move on the better move and what it does, every condition held.
\item The owner judged a day replayed with round 2 short of what he asked for: its scenes say what
  the better move is for, but a lesson is named rather than taught, and the line is played without
  saying why each move works, so it builds no pattern he can recognise. The audit's items do not
  ask for that, which is why the rule passed a version he did not want.
\item The next version reads every move of the line his move allowed and of the line the pick
  starts on the board (a check, a trade, a take-back, a threat, what the move is for), and its
  prompt teaches the lesson's method on his position before any move, plays his move's line to
  what it costs, gives each move of the better line its job, and ends on the cue that the idea is
  on. Measuring it needs items for those, an audit that each move's stated job matches the board,
  and the owner's own judgement of a replayed day.
\item False claims rose a little in round 2, inside the noise. A version that says more about each
  move has more to get wrong, so its round should hold false board claims at least as tightly as
  this rule does.
\end{itemize}

\section{Caveats}

\begin{itemize}[nosep]
\item \textbf{One run per config and day.} There is no second run of today's code for a noise floor;
  the same-script and decoy controls stand in for it.
\item \textbf{A model judges.} The judge and the audit are Sonnet 5.5 with the evidence, not the
  player; the controls show the judge neither invents differences nor rewards length.
\item \textbf{The days are past days.} The writer reads the same data the routine read, but the
  lessons' server serves today's lessons.
\item \textbf{Round 2 was designed after round 1.} Its sentence answers a gap round 1's audits of
  the same days showed, held-out days included, so round 2 is not an independent test: a new set of
  days would be.
\end{itemize}

\end{document}
